% Compile with: pdflatex project_presentation.tex (twice)
% Or: tectonic project_presentation.tex
% Self-contained source; no external images, fonts, or bibliography.
% Four-minute script and timing appear before each frame.
\documentclass[aspectratio=169,11pt]{beamer}
\usepackage[T1]{fontenc}
\usepackage{lmodern}
\usepackage{booktabs}
\definecolor{Ink}{HTML}{142C43}
\definecolor{Teal}{HTML}{007F82}
\definecolor{Muted}{HTML}{526575}
\definecolor{Paper}{HTML}{F5F8FA}
\definecolor{Mint}{HTML}{62DED0}
\definecolor{Gold}{HTML}{F3C66B}
\definecolor{RedChannel}{HTML}{C64A56}
\definecolor{GreenChannel}{HTML}{23826A}
\definecolor{BlueChannel}{HTML}{376DC6}
\setbeamercolor{normal text}{fg=Ink,bg=Paper}
\setbeamercolor{background canvas}{bg=Paper}
\setbeamercolor{palette primary}{fg=white,bg=Ink}
\setbeamercolor{frametitle}{fg=Ink}
\setbeamercolor{structure}{fg=Teal}
\setbeamercolor{alerted text}{fg=Teal}
\setbeamerfont{frametitle}{size=\LARGE,series=\bfseries}
\setbeamerfont{title}{size=\LARGE,series=\bfseries}
\setbeamertemplate{navigation symbols}{}
\setbeamertemplate{itemize item}{\small\textbullet}
\setbeamertemplate{footline}{%
\begin{beamercolorbox}[wd=\paperwidth,ht=4mm,dp=2mm,leftskip=9mm,rightskip=9mm]{palette primary}
\scriptsize\textbf{BINDU}\hspace{2mm}{\color{Mint}An image workspace}\hfill\insertframenumber\enspace/\enspace\inserttotalframenumber
\end{beamercolorbox}}
\setbeamersize{text margin left=9mm,text margin right=9mm}
\setbeamertemplate{frametitle}{\vspace{4mm}\insertframetitle\par\vspace{1mm}{\color{Teal}\rule{\textwidth}{0.7pt}}\vspace{1mm}}
\setbeamercolor{math text}{fg=Teal}
\setbeamercolor{description item}{fg=Teal}
\setbeamertemplate{enumerate item}{\color{Teal}\bfseries\insertenumlabel.}
\newcommand{\topic}[1]{{\color{Teal}\bfseries #1}}
\newcommand{\muted}[1]{{\color{Muted}#1}}
\title{Bindu: an image workspace}
\hypersetup{pdftitle={Bindu: an image workspace},pdfsubject={Four-minute project presentation}}
\begin{document}

% SLIDE 1: 0:00--0:10 (10 seconds)
% Our project is Bindu: an image workspace. We explore
% image channels, reconstruction, color spaces, Fourier compression and filters.
{
\setbeamercolor{background canvas}{bg=Ink}
\setbeamercolor{normal text}{fg=white,bg=Ink}
\usebeamercolor[fg]{normal text}
\begin{frame}[plain]
{\color{Mint}\small IMAGE PROCESSING PROJECT}\par\vspace{7mm}
{\fontsize{40}{44}\selectfont\bfseries Bindu:}\par\vspace{1mm}
{\fontsize{26}{30}\selectfont an image workspace}\par
\vspace{7mm}
{\color{Mint}\rule{22mm}{1.2pt}}\par\vspace{4mm}
{\large Channel analysis, reconstruction\\and image processing}\par
\vspace{7mm}
{\color{Gold}\small Python\quad /\quad Streamlit\quad /\quad NumPy\quad /\quad OpenCV}
\end{frame}
}

% SLIDE 2: 0:10--0:25 (15 seconds)
% The application has five main tools in this presentation. We analyze RGB
% channels, reconstruct selected channels, compare color spaces, compress with
% Fourier energy pruning, and apply smoothing or sharpening filters.
\begin{frame}{Feature overview}
\renewcommand{\arraystretch}{1.65}
\begin{tabular}{@{}lp{40mm}p{85mm}@{}}
\topic{01}&\textbf{Channel analyzer}&RGB previews, histograms and FFT spectra\\
\topic{02}&\textbf{Reconstruction}&Recover selected color channels\\
\topic{03}&\textbf{Color spaces}&Compare RGB and YCbCr\\
\topic{04}&\textbf{Fourier compression}&Prune coefficients and decode archives\\
\topic{05}&\textbf{Image filters}&Smooth or sharpen image details\\
\end{tabular}
\vspace{5mm}
\small\muted{Upload an image or choose a sample, then explore the processing results.}
\end{frame}

% SLIDE 3: 0:25--1:00 (35 seconds)
% Our project lets users explore how images work through an interactive
% Streamlit application. The channel analyzer starts by separating an image
% into red, green and blue components. A grayscale preview shows the intensity
% of each component, while a colored preview shows its contribution to the
% original image. Histograms count pixels at each intensity, helping us compare
% how the channels use their available range. The two-dimensional FFT gives a
% different view: it describes spatial frequencies rather than pixel positions.
% We display log-magnitude spectra so weaker frequency components remain visible.
% Together, these views connect the image's colors with its intensity and
% frequency information before we modify or reconstruct it.
\begin{frame}{Channel analyzer}
\textbf{Inspect each RGB component in the spatial and frequency domains}
\vspace{5mm}
\begin{itemize}\setlength\itemsep{4mm}
\item \textbf{Channel previews:} isolate red, green and blue contributions.
\item \textbf{Histograms:} compare pixel intensity distributions.
\item \textbf{2-D FFT spectra:} inspect spatial frequency content.
\end{itemize}
\vspace{4mm}
\[F_c(u,v)=\mathcal F\{I_c(x,y)\},\qquad c\in\{R,G,B\}\]
\vspace{2mm}
\muted{Log-magnitude display makes weaker frequency components easier to see.}
\end{frame}

% SLIDE 4: 1:00--1:35 (35 seconds)
% Partial reconstruction asks what happens when we keep only some color
% channels. The app computes the full complex FFT of each selected channel,
% then applies the inverse FFT to recover its spatial values. Both magnitude
% and phase matter in this reconstruction. Channels that are not selected
% become zero. For example, choosing red and green preserves those components
% and removes blue, changing the overall color balance. Users can select one
% channel or any pair, compare the result with the original and download it.
% This operation removes entire color channels. Within each selected channel,
% it retains the full frequency spectrum, so it differs from the coefficient
% pruning used later in Fourier compression.
\begin{frame}{Partial image reconstruction}
\begin{columns}[T,onlytextwidth]
\column{0.49\textwidth}
\topic{Keep one channel or a pair}
\begin{itemize}\setlength\itemsep{4mm}
\item Choose R, G, B, R+G, R+B or G+B
\item Apply the inverse FFT to selected channels
\item Set omitted channels to zero
\end{itemize}
\column{0.46\textwidth}
\topic{Full complex spectra}
\vspace{4mm}
\[\widehat I_c=\begin{cases}
\mathcal F^{-1}(F_c),&c\text{ selected},\\
0,&c\text{ omitted}.
\end{cases}\]
\vspace{3mm}
\muted{Example: R+G preserves red and green values and removes blue.}
\end{columns}
\vspace{6mm}
\alert{Selected channels retain both magnitude and phase.}
\end{frame}

% SLIDE 5: 1:35--2:10 (35 seconds)
% The color-space comparison shows the same image in RGB and full-range
% YCbCr. RGB expresses each pixel through red, green and blue intensities.
% YCbCr instead separates luma, called Y, from two color-difference channels,
% Cb and Cr. The displayed formula shows that green contributes the most to
% the luma calculation used here. A neutral color has chroma values of 128.
% The app compares channel previews, histograms and frequency spectra. Shared
% axes and a shared spectrum scale support fair visual comparison. By default,
% it removes channel means before plotting spectra so constant offsets do not
% dominate. This feature demonstrates the representation change without
% performing compression or chroma subsampling.
\begin{frame}{Color spaces: RGB and YCbCr}
\begin{columns}[T,onlytextwidth]
\column{0.45\textwidth}
{\Large\bfseries\color{RedChannel}R\color{GreenChannel}G\color{BlueChannel}B}
\par\vspace{3mm} Red, green and blue intensities describe each pixel.
\par\vspace{6mm}
{\Large\topic{YCbCr}}
\par\vspace{3mm} Y represents luma. Cb and Cr represent color differences.
\column{0.50\textwidth}
\topic{Full-range BT.601 conversion}
\vspace{3mm}
\[Y=0.299R+0.587G+0.114B\]
\muted{Neutral chroma: Cb = Cr = 128.}
\par\vspace{5mm} Compare grayscale previews, histograms and FFT spectra.
\end{columns}
\vspace{6mm}
\small\muted{Shared plotting scales support comparison. Mean removal is on by default for spectra. This tool does not compress or subsample the image.}
\end{frame}

% SLIDE 6: 2:10--3:10 (60 seconds)
% The lossy method transforms each RGB channel with a two-dimensional FFT.
% It groups conjugate frequency pairs and ranks them by squared magnitude,
% which represents energy. It always retains the DC component, or channel
% mean, then keeps enough of the strongest pairs to meet the chosen energy
% target. The archive stores indices and complex coefficients. Decoding rebuilds
% the spectrum and applies the inverse FFT without needing the original image.
% The table is a measurement using one resized repository sample. Raising the
% target from 99 to 99.9 percent improves PSNR but greatly increases file size.
% Energy retention is therefore not a compression ratio, and the archive can
% be larger than a PNG. These figures illustrate one image, not a benchmark.
\begin{frame}{Fourier compression}
\begin{columns}[T,onlytextwidth]
\column{0.49\textwidth}
\begin{enumerate}\setlength\itemsep{2mm}
\item Compute a 2-D FFT per channel
\item Rank conjugate pairs by energy
\item Keep DC and the strongest pairs
\item Store indices and coefficients
\end{enumerate}
\vspace{2mm}
\[\frac{\sum_{k\in K}|F(k)|^2}{\sum_k|F(k)|^2}\geq \frac{p}{100}\]
\muted{Decode with an inverse FFT.}
\column{0.46\textwidth}
\topic{Measured quality--size tradeoff}\par\vspace{3mm}
\small
\renewcommand{\arraystretch}{1.35}
\begin{tabular}{@{}rrr@{}}
\toprule
\textbf{Energy} & \textbf{Archive} & \textbf{PSNR}\\
\midrule
95\% & 2.7 KiB & 17.09 dB\\
99\% & 15.0 KiB & 24.09 dB\\
99.9\% & 193.8 KiB & 34.20 dB\\
\bottomrule
\end{tabular}
\par\vspace{4mm}
\muted{Higher PSNR means less pixel error.\\Energy retention is not a size ratio.}
\end{columns}
\vspace{5mm}
{\scriptsize\muted{Measured with the project codec on inputs/test.jpg resized to 256 $\times$ 366 RGB pixels. Raw size: 274.5 KiB. One example, not a general benchmark.}}
\end{frame}

% SLIDE 7: 3:10--3:50 (40 seconds)
% Finally, the image filters let us modify local image structure. Gaussian
% blur uses weighted neighbors, while box blur uses their average. Both smooth
% the image, with stronger settings also removing fine detail. For sharpening,
% the app provides Laplacian sharpening and unsharp masking. Unsharp masking
% subtracts a blurred image to isolate detail, then adds a scaled version of
% that detail back to the original. Users can compare filtering only red,
% green or blue against filtering all channels. Applying the same filter
% independently to every channel matches its whole-image counterpart within
% floating-point tolerance. These controls show the tradeoff between smoothing
% detail and emphasizing edges. Excessive sharpening can amplify noise or halos.
\begin{frame}{Image filters}
\begin{columns}[T,onlytextwidth]
\column{0.46\textwidth}
{\Large\topic{Smoothing}}
\begin{itemize}\setlength\itemsep{3mm}
\item Gaussian blur: weighted neighbors
\item Box blur: local averaging
\end{itemize}
\par\vspace{5mm}\muted{Larger blur settings smooth more detail.}
\column{0.49\textwidth}
{\Large\topic{Sharpening}}
\begin{itemize}\setlength\itemsep{3mm}
\item Laplacian: emphasize local changes
\item Unsharp mask: add detail back
\end{itemize}
\[I_{\mathrm{sharp}}=I+a\bigl(I-G_\sigma*I\bigr)\]
\small\muted{$G_\sigma*I$: Gaussian-blurred image.\\$a$: sharpening amount.}
\end{columns}
\vspace{5mm}
\alert{Compare R-only, G-only, B-only and all-channel filtering.}
\par\vspace{3mm}
\small\muted{Strong sharpening can amplify noise and produce halos.}
\end{frame}
% SLIDE 8: 3:50--4:00 (10 seconds)
% These tools connect image-processing concepts with visible results and
% demonstrate how channel selection, compression and filtering affect images.
% Thank you. I welcome your questions.
{
\setbeamercolor{background canvas}{bg=Ink}
\setbeamercolor{normal text}{fg=white,bg=Ink}
\usebeamercolor[fg]{normal text}
\begin{frame}[plain]
{\color{Mint}\small BINDU: AN IMAGE WORKSPACE}\par\vspace{10mm}
{\fontsize{36}{40}\selectfont\bfseries Thank you}\par
\vspace{5mm}
{\color{Gold}\Large Questions and discussion}\par
\vspace{8mm}
{\color{Mint}\rule{22mm}{1.2pt}}\par\vspace{4mm}
\large Exploring how channel selection, compression\\and filtering affect images.
\end{frame}
}
\end{document}
